(1) 探测记录 $I=|O+R_0e^{ik_rx}|^2=|O|^2+R_0^2+R_0Oe^{-ik_rx}+R_0O^*e^{ik_rx}$。直流包含两个自强度项，而 $\pm k_r$ 附近承载物光及其复共轭的两个互为镜像侧带。

(2) 若 $\widehat O(k)$ 仅支持 $[-B,B]$，交叉项占据约 $[-k_r-B,-k_r+B]$ 和 $[k_r-B,k_r+B]$。$|O|^2$ 的频谱支持最多 $[-2B,2B]$，因此充分条件可取 $|k_r|>3B$。选取例如 $-k_r$ 侧带并平移到零频，除以已知 $R_0$ 即可恢复 $O(x)$；另一侧带重建的是 $O^*$，即共轭双像。

(3) 局部令 $|O|=A$，强度在空间载频下 $I(x)=A^2+R_0^2+2AR_0\cos(\phi-k_rx)$，所以可见度为 $\mathcal V=(I_{\max}-I_{\min})/(I_{\max}+I_{\min})=2AR_0/(A^2+R_0^2)\le1$。只有 $A=R_0$ 时达到1；参考过强或过弱均降低调制度。

(4) 若参考相位不是已知的 $k_rx$，重建物光只确定到额外未知的参考相位因子。像素间距 $\Delta x$ 需使最高载频及侧带满足Nyquist限制（最大采样频率 $\pi/\Delta x$），否则别名重叠。物光与参考光缺乏相干性时交叉项时间平均消失。双像本质来自实强度记录同时含 $O$、$O^*$，离轴频域滤波可选择其中一份。
